\documentclass[10pt,a4paper]{amsart}
\usepackage[margin=19mm]{geometry}
\usepackage{amsmath,amssymb,mathtools}
\usepackage[scr=rsfs,cal=euler]{mathalfa}
\usepackage{xfrac}
\usepackage[unicode,hidelinks,bookmarks=false]{hyperref}
\newcommand{\ie}{i.e.\ }
\newcommand{\cf}{cf.\ }
\newcommand{\resp}{respectively\ }
\newcommand{\Subsection}{\subsection}
\providecommand{\nobreakdash}{\nobreak\hbox{-}}
\setlength{\emergencystretch}{2em}
\multlinegap0pt
\usepackage{bm}
\usepackage{bbm}
\usepackage{dsfont}
\usepackage{xspace}
\usepackage{cancel}
\usepackage{mathtools}
\usepackage{amsthm}
\makeatletter
\@ifundefined{thm}{\newtheorem{thm}{Thm}}{}
\@ifundefined{lem}{\newtheorem{lem}[thm]{Lemma}}{}
\@ifundefined{rem}{\newtheorem{rem}[thm]{Remark}}{}
\makeatother
\newcommand{\R}{\mathbb R}
\newcommand{\Z}{\mathbb Z}
\newcommand{\ads}{\mathbf{AdS}}
\newcommand{\bbe}{\mathbf{e}}
\newcommand{\bbx}{\mathbf{x}}
\newcommand{\bby}{\mathbf{y}}
\title[Counting Lattice Points in Generalized Permutohedra From A t]{Counting Lattice Points in Generalized Permutohedra From A to B}
\author{Warut Thawinrak}
\date{Source: arXiv:2512.01332v2 (CC BY 4.0)}
\begin{document}
\maketitle

\noindent\emph{Source and licence.} Counting Lattice Points in Generalized Permutohedra From A to B. Version arXiv:2512.01332v2 (CC BY 4.0), distributed under the Creative Commons Attribution 4.0 International licence, as recorded in the arXiv licence metadata for this exact version and shown on the abstract page https://arxiv.org/abs/2512.01332v2. Full rights notice: licenses/arxiv-2512.01332-CC-BY-4.0.txt.\par\smallskip

\noindent\textit{Editorial scope.} Counted unit: the Lem whose source label is \texttt{lem: no-common-interior} in the article (source0.tex lines 695--696, source label \texttt{lem: no-common-interior}), delivered together with the article's own complete original proof of it (source0.tex lines 698--712, one \texttt{proof} environment of 15 lines). No mathematics is written, repaired or re-derived: statement and proof are the article's own text line for line. Required context retained uncounted: rem: at-most-n-left-vertices (lines 691--693). Every definition, display and label of those lines is retained unchanged. Omitted from this package and not redistributed: 1--690, 694--694, 697--697, 713--870. Nothing inside a retained excerpt is omitted or altered: every retained body is delivered exactly as the article's lines are written, so no omission receipt of any kind accompanies this package. Citations: the retained excerpts carry no citation command at all, so no bibliography entry of the article is transcribed and no reference is redistributed. Reference adaptation: none was needed and none was made; every reference inside the delivered unit resolves inside it, so no unresolved reference or citation is produced. Printed slips kept literal: no printed word, symbol, hypothesis or number of the counted statement or of its proof was altered. Layout and class: the article's own class is \texttt{article} with packages including geometry, inputenc, amsfonts, amsmath, amsthm, amssymb, graphicx, tikz, enumitem, hyperref, xcolor, mathtools, relsize, caption; the wrapper is the lane's pinned \texttt{amsart} editorial wrapper at 10pt on A4, which restates the theorem-like environments the retained text needs and adds the article's own macro definitions that the retained text needs (\texttt{\textbackslash R}, \texttt{\textbackslash Z}, \texttt{\textbackslash ads}, \texttt{\textbackslash bbe}, \texttt{\textbackslash bbx}, \texttt{\textbackslash bby}), with extra wrapper packages (bm, bbm, dsfont, xspace, cancel, mathtools). The delivered numbering is the unit's own. Assets: no retained span contains a figure environment, a graphics-inclusion command, a PDF-page inclusion, a table or a TikZ picture; no image file is copied into this package, the package declares no asset dependency and no image file is redistributed with it. Licence: version arXiv:2512.01332v2 of this article is distributed under the Creative Commons Attribution 4.0 International licence, as recorded in the arXiv licence metadata for this exact version and shown on the abstract page https://arxiv.org/abs/2512.01332v2. A full rights notice is delivered at licenses/arxiv-2512.01332-CC-BY-4.0.txt. Characters: every retained excerpt is pure ASCII, so the delivered engine, XeLaTeX with its default Latin Modern text font, reports no missing character.
\begin{rem}\label{rem: at-most-n-left-vertices}
    If we assume further that $P_T$ is $n$-dimensional ($G(P_T)$ is connected), then $H$ is a spanning tree of $G(P_T)$ on $m$ left and $n+1$ right vertices. Consequently, we have that $\hat{H}$ is a bipartite graph on $n+1$ right vertices in which every left vertex has degree at least two. Moreover, $\hat{H}$ has at most $n$ left vertices. The tree $\hat{H}$ has exactly $n$ left vertices if and only if every left vertex of $\hat{H}$ has degree two.
\end{rem} 

\begin{lem}\label{lem: no-common-interior} Let $P = \Delta^0_{S_1} + \cdots + \Delta^0_{S_m}$ where $S_i \in \ads$ for all $i \in [m]$. Suppose that $T \in \ads_n$ is an admissible subset in which $P_T$ is $n$-dimensional. If $\Pi$ is a fine mixed cell of $P_T$, then $\Pi$ and $\bby + \Pi$ have no common interior for all vectors $\bby \in \Z^n\backslash\{\mathbf{0}\}.$
\end{lem} 

\begin{proof}
    Suppose that $\Pi$ is a fine mixed cell of the form $\Pi = \Delta_{I_1} + \cdots + \Delta_{I_m}$ where  $I_i \subseteq \{0\}\cup S_i\cap T$ for all $i \in [m]$. We claim that $\Pi$ lies in an integral translation of a \emph{fundamental parallelepiped}. It is well-known that any fundamental parallelepiped has no common interior with its integral translation. Thus, the conclusion of the lemma follows from this claim. Hence, it only remains to prove the claim.
    
    Let $K = \{I_i\ | \ i \in [m] \text{ and } |I_i| \geq 2\}$. Then, 
    \[\hat{\Pi} := \sum_{I \in K}\Delta_I.\]
    For each $I \in K$, let us write $I= \{j_1, \dots, j_{|I|}\ | \ |j_i| < |j_{i+1}|, \text{ for all } 1 \leq i < |I| \}$, and define the corresponding set of intervals 
    $\binom{I}{2}^* := \{[\bbe_{j_1}, \bbe_{j_2}], [\mathbf{0}, \bbe_{j_3} - \bbe_{j_1}], \dots, [\mathbf{0},\bbe_{j_{|I|}} - \bbe_{j_1}] \subset \R^n\}$
    where we denote by $[\bbx, \bby]$ the line segment connecting $\bbx$ and $\bby$ (the interval from $\bbx$ to $\bby$), and set $\bbe_0 = \mathrm{0} \in \R^n$. It is easy to see that every vertex of the simplex $\Delta_I$ lies in the zonotope
    \[\Diamond_I := \sum_{\Delta \in \binom{I}{2}^*}\Delta\]
    spanned by the intervals in $\binom{I}{2}^*.$ This implies that $\Delta_I \subseteq \Diamond_I$ for all $I \in K$. In addition, let
    \[B_K := \bigcup_{I \in K} \binom{I}{2}^*\]
    be the set of intervals from all of $\binom{I}{2}^*$ for $I \in K$. Thus, $\hat{\Pi}$ lies in the zonotope
    \[\Diamond_K := \sum_{\Delta \in B_K}\Delta = \sum_{I \in K}\Diamond_I.\]
    Since $P_T$ is $n$-dimensional, by Remark \ref{rem: at-most-n-left-vertices}, $\hat{H}$ is a tree and a bipartite graph with $n+1$ right vertices. In addition, because $|\binom{I}{2}^*| = |I|-1$, it follows that $$\sum_{I \in K}\Big|\binom{I}{2}^*\Big| = n.$$ Hence, the polytope $\Diamond_K$ is an integral translation of a fundamental parallelepiped in $\R^n$. Since $\hat{\Pi} \subseteq \Diamond_K$ and $\Pi$ is an integral translation of $\hat{\Pi}$, it follows that $\Pi$ lies in an integral translation of a fundamental parallelepiped as claimed.
\end{proof}

\end{document}
